% Candidate colour pairs for potassium and rubidium if green and violet leave the palette (study).
\documentclass[10pt,border=1mm]{standalone}
\usepackage[T1]{fontenc}
\usepackage{figurestyle}
\input{primitives.tex}
\input{molecules.tex}
\input{numbers.tex}
% Nicker poster colours as read off the author's chart (approximate hex, by eye)
\definecolor{NkFrenchGray}{HTML}{9A9FA3}\definecolor{NkBlack}{HTML}{1E1A1C}
\definecolor{NkCarmine}{HTML}{8E2A2E}\definecolor{NkScarlet}{HTML}{B33A2C}\definecolor{NkSienna}{HTML}{8C5A48}
\definecolor{NkMauve}{HTML}{4E3A8E}\definecolor{NkCobaltViolet}{HTML}{5F3F95}\definecolor{NkPink}{HTML}{EC7FB0}
\definecolor{NkMagenta}{HTML}{A3338C}\definecolor{NkOpera}{HTML}{C9435F}\definecolor{NkOrange}{HTML}{E8683A}
\definecolor{NkPrussian}{HTML}{2E3A6B}\definecolor{NkCobalt}{HTML}{2E5CB8}\definecolor{NkCerulean}{HTML}{2A79B5}
\definecolor{NkFrenchBlue}{HTML}{2F5A9E}\definecolor{NkLightBlue}{HTML}{2FA9D9}\definecolor{NkLemon}{HTML}{F5E73A}
\definecolor{NkViridian}{HTML}{0F7A58}\definecolor{NkImperial}{HTML}{2C9E4A}\definecolor{NkLightGreen}{HTML}{3AA85A}
\definecolor{NkChromeGreen}{HTML}{8FBF3A}\definecolor{NkOcher}{HTML}{BF9433}\definecolor{NkChromeYellow}{HTML}{F2C23A}

\definecolor{ColTeal}{HTML}{4F9C97}
\begin{document}
\begin{tikzpicture}[plate]
\path[use as bounding box] (0,0) rectangle (15,17.5);
\newcommand\row[4]{% y, K colour, Rb colour, caption
  \begin{scope}[shift={(0,#1)}]
    \colorlet{ColPotassium}{#2}\colorlet{ColRubidium}{#3}
    \begin{scope}\molsolid\mollabelsfalse
      \pic[scale=.62] at (2.2,1.2) {mol3d-krb-in}; \pic[scale=.62] at (5.6,1.2) {mol3d-krb-complex}; \pic[scale=.62] at (9.0,1.2) {mol3d-krb-out};
      \pic[scale=.55] at (11.9,1.2) {mol3d-mla-ref}; \pic[scale=.55] at (13.9,1.2) {mol3d-water-X};
    \end{scope}
    \path[fill=ColTeal!38] (.3,2.45) rectangle (.9,2.75); \path[fill=ColGold!55] (1.0,2.45) rectangle (1.6,2.75);
    \node[sub,anchor=west] at (1.8,2.6) {#4};
  \end{scope}}
\row{14.3}{NkOcher}{ColTeal}{1. K ocher, Rb teal}
\row{11.5}{ColTeal}{NkOcher}{2. K teal, Rb ocher}
\row{8.7}{NkSienna}{NkFrenchGray}{3. K burnt sienna, Rb French gray}
\row{5.9}{NkOcher}{NkSienna}{4. K ocher, Rb burnt sienna}
\row{3.1}{NkFrenchGray}{NkPrussian}{5. K French gray, Rb Prussian blue (blue is then no longer nitrogen's alone)}
\row{.3}{ColGold}{PlateInk!55}{6. K gold, Rb mid grey}
\node[sub,anchor=north west] at (.3,17.4) {each row: reactants, complex, products; then methylamine and water for the suite context; the two swatches are the teal sheet tint and the gold retained tint};
\end{tikzpicture}
\end{document}
